\section{Emergent Abilities: ¿salto de capacidad o artefacto de medición?}

La clase empieza por recordar la definición clásica de 2022 sobre las capacidades emergentes de los modelos grandes: una capacidad está cerca del nivel aleatorio en los modelos pequeños, aparece con claridad solo a partir de cierta escala y no puede extrapolarse directamente de la tendencia de los modelos pequeños. El ponente añade de inmediato la controversia sobre la medición: las métricas discretas, el prompt y la definición de la tarea pueden cambiar la forma de la curva. Por ello, la pregunta correcta no es la disyuntiva «¿la emergencia es real o falsa?», sino qué conductas, con qué medición y por qué se manifiestan como un umbral.

\subsection{Definición clásica y critical threshold}

La primera diapositiva de título presenta la emergencia como un objeto de estudio en sí mismo; luego, la diapositiva de definición plantea tres condiciones: los modelos grandes la tienen y los pequeños no, no se puede extrapolar directamente a partir de los modelos pequeños y el desempeño es casi aleatorio antes del critical threshold y sube con rapidez después. Esta definición se centra en la observed behavior y no ofrece una explicación del mecanismo interno. Esta sección fija primero esa definición histórica; más adelante se ponen a prueba los límites de su evidencia con la controversia sobre el prompt y la metric.

\lecturefigure{slide-02-emergent-title.jpg}{Emergent Abilities of Large Language Models}{Video de la clase de Stanford 00:00:33}

\begin{knowledgebox}{Lectura de la figura: la diapositiva de título establece el objeto de estudio}
Aquí, la ability puede ser modular arithmetic, word unscrambling, question answering o una prompting behavior. No se trata de que una sola neuron «aprenda una habilidad» de golpe, sino de un patrón de desempeño en la tarea que observa la evaluation pipeline; el modelo, el prompt, el scoring y la task distribution forman parte de la definición del fenómeno.
\end{knowledgebox}

\lecturefigure{slide-03-emergent-definition.jpg}{Definición de capacidad emergente que adopta la clase}{Video de la clase de Stanford 00:00:39}

\begin{importantbox}{Intuición matemática de la definición por umbral}
Sea $s$ la escala y $m(s)$ la puntuación en la tarea. La narrativa clásica de la emergencia sostiene que existe un $s_c$ tal que
\[
m(s)\approx m_{\text{random}}\quad(s<s_c),\qquad
m(s)\gg m_{\text{random}}\quad(s\ge s_c).
\]
$s_c$ es la critical scale observada y $m_{\text{random}}$ es la línea base aleatoria. Esta expresión describe la forma de la curva; no demuestra que en el espacio de parámetros ocurra una verdadera transición de fase física.
\end{importantbox}

\teachervoice{El ponente usa la phase transition como analogía del umbral de escala, pero también dice que los investigadores todavía carecen de una explicación suficiente. Lo «impredecible» en la clase significa que una extrapolación ingenua a partir de las curvas medidas en modelos pequeños falla, no que ninguna capacidad pueda estudiarse ni predecirse.}

\subsection{Cómo leer las curvas few-shot}

La figura con varias tareas pone los training FLOPs en el eje horizontal y la performance de la tarea en el eje vertical; modelos como LaMDA, GPT-3, Gopher, Chinchilla y PaLM muestran umbrales distintos en modular arithmetic, IPA transliteration, word unscrambling y Persian QA. Al leer la figura no basta con ver que «el modelo más grande obtiene la puntuación más alta»: también hay que revisar si el tramo de los modelos pequeños es plano, si la tarea usa exact match y si las distintas familias de modelos siguen la misma tendencia.

\lecturefigure{slide-04-few-shot-emergence.jpg}{Curvas de varias tareas few-shot en función de los training FLOPs}{Video de la clase de Stanford 00:01:18}

\begin{knowledgebox}{Lectura de la figura: primero se confirman los ejes y después se juzga el salto}
El eje horizontal es la escala logarítmica del cómputo de entrenamiento; el eje vertical varía según la tarea. Algunas gráficas suben con claridad alrededor de $10^{22}$--$10^{23}$ FLOPs, pero hay pocos puntos y las arquitecturas de los modelos difieren. Si la metric es de todo o nada, una mejora continua a nivel de token puede comprimirse en un crossing repentino; por eso, el umbral puede reflejar tanto un cambio real de capacidad como una scoring nonlinearity.
\end{knowledgebox}

\begin{warningbox}{No se debe deducir una transición de fase necesaria a partir de unos pocos puntos discretos}
La escala del modelo, los datos, la arquitectura y la recipe de entrenamiento cambian al mismo tiempo; los puntos de la curva no provienen de un controlled experiment que modifique una sola variable independiente. Las gráficas de emergencia respaldan que «la extrapolación ingenua desde modelos pequeños no es confiable», pero no respaldan por sí mismas que «una representación interna aparece de repente en cierto punto».
\end{warningbox}

\subsection{La prompting strategy cambia la capacidad que se observa}

El siguiente grupo de figuras muestra el augmented prompting, que incluye scratchpad, instruction following, suma de ocho dígitos y calibration. El ejemplo más representativo es GSM8K: con prompting ordinario la mejora con la escala es limitada, mientras que Chain-of-Thought aumenta el desempeño de forma notable en los modelos grandes. Aquí la capacidad proviene tanto de los parámetros como de exponer el proceso de cálculo en los token de salida.

\lecturefigure{slide-05-augmented-prompting.jpg}{Cambios de escala y capacidad con augmented prompting}{Video de la clase de Stanford 00:01:54}

\begin{knowledgebox}{Lectura de la figura: la capacidad del modelo no puede separarse del elicitation method}
Un mismo checkpoint puede producir curvas completamente distintas con prompts diferentes. Si CoT hace que el modelo grande convierta un cálculo implícito de varios pasos en una secuencia explícita, la «emergencia» observada puede ser la combinación de una latent capability con un elicitation threshold. Por eso, un informe de evaluación debe declarar a la vez el model, los prompt examples, el decoding y la metric.
\end{knowledgebox}

\begin{importantbox}{El desempeño observado es una función compuesta}
Escrita de forma más completa, la puntuación de la tarea es
\[
m=G(\theta, p, d, \mu),
\]
donde $\theta$ representa los parámetros y el entrenamiento del modelo, $p$ la estrategia de prompt/inferencia, $d$ el decoding y $\mu$ la metric. Si solo se grafica $m$ contra la escala de parámetros, la no linealidad de los otros tres términos se proyecta junto con ella en el «efecto de escala».
\end{importantbox}

\subsection{El catálogo de capacidades y la escala son solo un índice}

La clase conserva la tabla grande del artículo, que enumera distintas emergent abilities con sus training FLOPs, número de parámetros, modelos y referencias. Su valor está en permitir a los investigadores encontrar patrones entre tareas, no en asignar a cada capacidad un «número mínimo de parámetros» fijo y permanente. Por ello, esta sección trata la tabla como un índice de experimentos y se concentra en revisar la tarea, la familia de modelos y la configuración de evaluación, en lugar de memorizar las cifras umbral.

\lecturefigure{slide-06-emergent-abilities-table.jpg}{Capacidades emergentes y escalas de aparición recopiladas en el artículo}{Video de la clase de Stanford 00:02:33}

\begin{knowledgebox}{Lectura de la figura: la scale de la tabla no es un umbral universal}
Una misma tarea desplaza su umbral con distintos tokenizer, data mixture, instruction tuning y prompt; además, cada familia de modelos aprovecha sus parámetros de manera diferente. Primero se mira la task y la reference, y después los FLOPs/parameters. La columna de escala funciona más como índice histórico de experimentos que como una constante que pueda consultarse directamente al diseñar un modelo nuevo.
\end{knowledgebox}

\subsection{Por qué ocurre la emergencia sigue sin resolverse}

La diapositiva de explicaciones dice de forma explícita que por ahora se sabe poco sobre las causas; las evaluation metrics utilizadas tampoco bastan para explicar el fenómeno, y se propone analizarlo con métricas continuas como la cross-entropy loss. Una metric continua puede revelar mejoras graduales que ya ocurren en los modelos pequeños, pero que todavía no cruzan el umbral de exact match. Esta sección pasa de «ver el salto» a «explicar el salto» y analiza por separado el measurement artifact y el mecanismo de la capacidad.

\lecturefigure{slide-07-emergence-explanations.jpg}{Incertidumbre sobre las causas de la emergencia y la evaluation metric}{Video de la clase de Stanford 00:02:45}

\begin{knowledgebox}{Lectura de la figura: la metric critique cambia la fuerza de la evidencia}
Si exact match solo otorga 1 punto cuando la respuesta completa es correcta, el paso del modelo de «casi correcto» a «correcto» produce un salto; la cross-entropy, en cambio, muestra cómo la masa de probabilidad se desplaza poco a poco. Una curva continua puede debilitar la impresión visual de una «transición de fase repentina», pero no descarta que una conducta compleja solo sea utilizable a partir de cierta escala práctica.
\end{knowledgebox}

\begin{warningbox}{La Mirage critique tampoco significa que «la capacidad no cambió»}
Que la métrica produzca una discontinuity no implica que la escala no tenga efecto; indica que «la capacidad aparece de repente» y «la puntuación cruza de repente una línea» deben distinguirse. En ingeniería sigue importando cuándo una capacidad alcanza un umbral utilizable; lo que no se puede es tomar ese umbral como la única evidencia del mecanismo.
\end{warningbox}

\teachervoice{El ponente no propone una única emergence theory, sino que sitúa la elección de la metric como una pregunta central de investigación. Este límite prudente importa más que memorizar umbrales: primero hay que preguntar cómo se mide el fenómeno y después qué ocurrió dentro del modelo.}

\subsection{Líneas de investigación Beyond scaling}

Aunque aumentar la escala pueda traer capacidades nuevas, la clase insiste en que no es la única variable. Nuevas arquitecturas, datos de mayor calidad, mejores training procedures, few-shot prompting más eficaz, la teoría y la interpretabilidad, y la computational linguistics pueden hacer que las capacidades aparezcan antes o sean más controlables.

\lecturefigure{slide-08-beyond-scaling.jpg}{Líneas de investigación sobre capacidades además de seguir aumentando la escala}{Video de la clase de Stanford 00:03:06}

\begin{knowledgebox}{Lectura de la figura: estas líneas actúan sobre capas distintas}
La architecture y el training modifican cómo se aprenden los parámetros; la data quality modifica la señal de supervisión; el prompting modifica la inference interface; la interpretabilidad y la theory modifican la capacidad de explicar y predecir; la computational linguistics aporta priors estructurales. Si todas se reducen a «cómputo efectivo», se pierde la posibilidad de intervenir en cada una.
\end{knowledgebox}

\subsection{El crecimiento de las capacidades va acompañado del crecimiento de los riesgos}

La Emergent risk slide reúne truthfulness, bias, toxicity, memorization y riesgos desconocidos. Un modelo más grande puede resolver mejor las tareas, pero también reproducir con más eficacia datos sensibles, adaptarse a prompts dañinos o amplificar sesgos en contextos ambiguos; capacidad y riesgo no comparten un mismo eje de seguridad monótono. Por ello, esta sección trata la evaluación de seguridad como una dimensión independiente, en lugar de inferir a partir de la puntuación promedio en las tareas que el riesgo ya disminuyó.

\lecturefigure{slide-09-emergent-risks.jpg}{El aumento de escala puede amplificar a la vez las capacidades y los riesgos sociales}{Video de la clase de Stanford 00:04:00}

\begin{knowledgebox}{Lectura de la figura: el riesgo no es un apéndice que se agrega después del entrenamiento}
La truthfulness es un problema de hechos y de calibración; el bias involucra diferencias entre grupos; la toxicity, salidas dañinas; la memorization, privacidad y derechos de autor. Cada uno requiere su propio dataset, attack y metric. Una puntuación global de helpfulness no demuestra que todos estos riesgos disminuyan al mismo tiempo.
\end{knowledgebox}

\teachervoice{El ponente advierte que los modelos futuros podrían presentar riesgos que aún no se han descubierto. Mencionar el unknown risk no busca sembrar alarma, sino mostrar que tanto la capability discovery como la safety discovery requieren red teaming y measurement continuos.}

\subsection{El general-purpose model cambia la estructura de la investigación y de las aplicaciones}

La Sociological changes slide no trata solo de la difusión tecnológica. Los modelos de propósito general llevan a la comunidad de «entrenar un modelo por tarea» a una shared foundation que luego atiende varios dominios mediante in-context prompting o una adaptación ligera; el Transformer también se extendió de NLP a escenarios como text-to-image, biology y robotics.

\lecturefigure{slide-10-sociological-changes.jpg}{Cambios en la comunidad y en las aplicaciones traídos por la emergencia y los modelos de propósito general}{Video de la clase de Stanford 00:04:45}

\begin{knowledgebox}{Lectura de la figura: un paradigma técnico redistribuye las barreras de entrada}
El shared model reduce el costo de desarrollar aplicaciones, pero concentra en unos pocos proveedores el poder de entrenamiento, las interfaces de evaluación y la infraestructura. Los investigadores pueden acceder más rápido a las capacidades, pero dependen más de data/updates que no pueden ver. Más adelante, BabyLM plantea justamente la pregunta inversa ante esta concentración de recursos.
\end{knowledgebox}

\subsection{Trabajo futuro y preguntas abiertas de la clase}

La Future-work slide reúne escala, arquitectura, entrenamiento, datos, prompting, capacidades de modelos pequeños, theory y computational linguistics; la diapositiva de discusión pregunta además si la escala tiene un límite, qué papel desempeñan los datos y la metric en la emergence, y si los specialized systems podrían superar al monolith que «mete todo en los parámetros».

\lecturefigure{slide-11-future-work.jpg}{Líneas de investigación futura sobre las emergent abilities}{Video de la clase de Stanford 00:06:00}

\begin{knowledgebox}{Lectura de la figura: el future work no es un único scaling roadmap}
La lista incluye a la vez capability, efficiency, interpretabilidad y linguistic structure. Lo que más vale conservar es la comparación entre varias rutas: una misma capacidad puede obtenerse con más parámetros, mejores datos, mejores prompts, herramientas externas o modelos especializados; el objetivo de investigación debe comparar costo y confiabilidad, no solo el peak score.
\end{knowledgebox}

\lecturefigure{slide-12-discussion-questions.jpg}{Preguntas abiertas sobre la emergence que el ponente deja a la clase}{Video de la clase de Stanford 00:06:57}

\begin{knowledgebox}{Lectura de la figura: las preguntas mismas marcan el límite de la evidencia}
Estas preguntas no tienen respuesta estándar: si seguirán apareciendo capacidades, si los rendimientos de la escala son decrecientes, si la metric y los data fabrican los umbrales, si los sistemas especializados son más confiables. Las notas no deben proyectar hacia atrás las opiniones de 2026 como conclusiones del ponente, sino usar estas preguntas como punto de conexión con las tres unidades siguientes: reasoning, BabyLM y agents.
\end{knowledgebox}

\teachervoice{El ponente invita a los estudiantes a interrumpir y discutir en cualquier momento, lo que muestra que esta parte no anuncia una ley de la emergence, sino que usa el artículo clásico para construir un espacio común de preguntas.}

\subsection{Resumen de la sección}

La definición clásica de emergencia describe puntuaciones de tarea con forma de umbral, pero la curva depende a la vez del modelo, el prompt, el decoding y la metric. La escala puede traer capacidades nuevas, pero también riesgos nuevos y concentración de recursos; mejores datos, arquitecturas, procedimientos de inferencia y sistemas especializados son rutas alternativas o complementarias.
